
\begin{titlepage}
    \makeatletter
    \AddToShipoutPictureBG*{
        \begin{tikzpicture}[overlay,remember picture]
            \draw [line width=1pt]
            ($ (current page.north west) + (1.95cm,-2.1cm) $)
            rectangle
            ($ (current page.south east) + (-1.95cm,1.75cm) $);
        \end{tikzpicture}

    }
    \makeatother
\section{自我总结}

\noindent
\textbf{一、思想政治表现}

\par
\indent
入学以来，本人始终坚持正确的政治方向，认真参加学校、学院及班级组织的各项学习活动，不断提高思想政治素养，树立正确的世界观、人生观和价值观。在学习和科研过程中，坚持实事求是、诚实守信，严格遵守学校各项规章制度，恪守学术道德和科研诚信，始终以严谨认真的态度对待课题研究、实验数据处理及论文撰写，自觉维护良好的学术风气。

\par
\indent
作为非全日制研究生，本人在完成本职工作的同时，坚持按要求完成课程学习和科研任务，积极参加课题组组会、学术交流及各项培养环节，主动与导师沟通研究进展，不断提升自身综合素质。面对工作、学习和生活之间的平衡问题，逐渐形成了较好的时间管理能力和自我规划能力，也更加珍惜每一次学习和科研机会。

\noindent\textbf{二、课程学习与专业能力提升}

\par
\indent
按照培养方案要求，本人认真完成了培养计划规定的各项课程学习，系统掌握了软件工程及智能优化相关专业知识，在算法设计、软件工程方法、分布式系统、机器学习等方面进一步夯实了理论基础。课程学习不仅帮助我完善了专业知识体系，也使我能够将课堂所学与课题研究相结合，为后续开展网络化分布式优化、多智能体协同优化及强化学习相关研究提供了理论支撑。

\par
\indent
除课程学习外，本人坚持阅读国内外相关文献，持续关注网络化分布式优化、多智能体协同进化、黑箱优化、约束优化及强化学习等方向的发展动态，不断拓宽专业视野。在文献调研过程中，逐步掌握了科研论文阅读、问题分析及研究思路凝练的方法，对本课题的研究背景、国内外研究现状及关键科学问题有了更加深入的理解。同时，通过参加相关学术会议和学术交流活动，进一步了解了本领域的发展趋势，也增强了开展科研工作的信心。

\noindent\textbf{三、科研训练与课题进展}

\par
\indent
在导师胡晓敏教授的指导下，本人围绕学位论文《基于强化学习惩罚因子的多智能体协同优化方法》持续开展研究工作。课题以网络化分布式优化（Network-based Distributed Optimization，NDO）为研究背景，针对现有多智能体协同优化算法中惩罚参数依赖人工设定、通信协商成本较高以及共享变量处理效率不足等问题，探索将强化学习引入协同优化过程，提高算法在复杂黑箱优化场景下的自适应能力与协同效率。

\par
\indent
围绕上述研究目标，本人完成了强化学习惩罚因子调控模块、共享变量重要性评估机制以及智能通信门控策略等关键模块的设计与实现，并在原有 MACPO 框架基础上构建了 RL-MACPO 算法。通过不断修改程序、分析实验结果和优化算法结构，逐步形成了较完整的算法框架，并完成了核心代码开发、实验平台搭建以及数据统计分析等工作。

\par
\indent
目前，课题已取得较为扎实的阶段性成果。在标准基准测试中，完成了 RL-MACPO 与原始 MACPO 的系统对比实验，并围绕算法关键模块开展了消融实验、强化学习策略对比实验以及通信策略分析。同时，在多个扩展测试函数及典型应用场景中，对算法性能进行了进一步验证，并持续完善相关基线算法复现工作，为后续论文分析提供更加充分的实验依据。

\par
\indent
经过这一阶段的科研训练，本人逐步完成了从文献调研、问题分析、算法设计、程序实现到实验验证和论文撰写的完整科研过程。通过不断发现问题、分析问题和解决问题，不仅提高了科研实践能力，也增强了独立开展科研工作的能力。目前，课题已由算法设计阶段逐步进入实验完善、理论分析和论文撰写阶段，为后续完成学位论文奠定了较好的基础。

\end{titlepage}

% ===== 第 2 页（续；边框与第 1 页相同，顶距对齐首页正文）=====
\begin{titlepage}
    \makeatletter
    \AddToShipoutPictureBG*{
        \begin{tikzpicture}[overlay,remember picture]
            \draw [line width=1pt]
            ($ (current page.north west) + (1.95cm,-2.1cm) $)
            rectangle
            ($ (current page.south east) + (-1.95cm,1.75cm) $);
        \end{tikzpicture}
    }
    \makeatother
    \vspace{2.15cm}

\par
\textbf{(续)}
\par

\textbf{四、主要收获与体会}

\par
\indent
经过硕士阶段的学习和科研训练，本人最大的收获不仅是专业知识的积累，更重要的是科研思维和解决问题能力的提升。刚接触课题时，对网络化分布式优化、多智能体协同进化以及强化学习等内容的理解还停留在理论层面。随着文献阅读的不断深入、实验的反复验证以及导师的耐心指导，我逐渐能够围绕实际问题分析算法设计思路，并尝试提出自己的改进方案。

\par
\indent
科研过程中，我深刻体会到，一个可靠的实验结果往往需要经过大量重复验证才能真正说明问题。从程序调试、参数选择，到实验复现、统计分析，每一个细节都会影响最终结论。正是在一次次修改程序、排查问题和不断优化实验的过程中，我逐渐养成了严谨求实的科研态度，也更加理解了科研工作需要耐心、坚持和不断积累。

\par
\indent
作为非全日制研究生，在兼顾工作与学习的过程中，我更加珍惜每一次科研训练和学习机会，也不断提高了时间管理和自我规划能力。虽然科研过程中遇到过程序调试困难、实验周期较长以及论文写作压力等问题，但每解决一个问题，都让我对科研工作有了更加深入的理解，也更加坚定了继续深入研究的信心。

\noindent\textbf{五、存在不足与改进措施}

\par
\indent
对照培养目标和学位论文要求，本人认识到目前仍存在一些不足。首先，部分扩展实验规模较大、运行周期较长，实验数据仍需进一步补充和完善，后续需要进一步提高实验组织和资源调度效率。其次，论文理论分析部分仍有进一步深化空间，对强化学习调控机制、通信策略以及算法性能提升原因的分析还需更加系统和深入。此外，在论文整体结构和学术表达方面仍需进一步打磨，使论文逻辑更加严谨、内容更加完整。

\par
\indent
下一阶段，本人将在导师指导下继续完善各项实验内容，进一步补充基线算法对比与扩展实验，加强算法机理分析，认真完成学位论文撰写工作，不断提高论文质量。同时，将继续保持严谨务实的科研态度，认真完成硕士阶段各项培养任务，努力以更加扎实的研究成果完成学位论文。

\end{titlepage}

\begin{titlepage}
    \makeatletter
    \AddToShipoutPictureBG*{
        \begin{tikzpicture}[overlay,remember picture]
            \draw [line width=1pt]
            ($ (current page.north west) + (1.95cm,-2.1cm) $)
            rectangle
            ($ (current page.south east) + (-1.95cm,1.75cm) $);
        \end{tikzpicture}

    }
    \makeatother
\section{自我总结}

\noindent
\textbf{一、思想政治表现}

\par
\indent
入学以来，本人始终坚持正确的政治方向，认真参加学校、学院及班级组织的各项学习活动，不断提高思想政治素养，树立正确的世界观、人生观和价值观。在学习和科研过程中，坚持实事求是、诚实守信，严格遵守学校各项规章制度，恪守学术道德和科研诚信，始终以严谨认真的态度对待课题研究、实验数据处理及论文撰写，自觉维护良好的学术风气。

\par
\indent
作为非全日制研究生，本人在完成本职工作的同时，坚持按要求完成课程学习和科研任务，积极参加课题组组会、学术交流及各项培养环节，主动与导师沟通研究进展，不断提升自身综合素质。面对工作、学习和生活之间的平衡问题，逐渐形成了较好的时间管理能力和自我规划能力，也更加珍惜每一次学习和科研机会。

\noindent\textbf{二、课程学习与专业能力提升}

\par
\indent
按照培养方案要求，本人认真完成了培养计划规定的各项课程学习，系统掌握了软件工程及智能优化相关专业知识，在算法设计、软件工程方法、分布式系统、机器学习等方面进一步夯实了理论基础。课程学习不仅帮助我完善了专业知识体系，也使我能够将课堂所学与课题研究相结合，为后续开展网络化分布式优化、多智能体协同优化及强化学习相关研究提供了理论支撑。

\par
\indent
除课程学习外，本人坚持阅读国内外相关文献，持续关注网络化分布式优化、多智能体协同进化、黑箱优化、约束优化及强化学习等方向的发展动态，不断拓宽专业视野。在文献调研过程中，逐步掌握了科研论文阅读、问题分析及研究思路凝练的方法，对本课题的研究背景、国内外研究现状及关键科学问题有了更加深入的理解。同时，通过参加相关学术会议和学术交流活动，进一步了解了本领域的发展趋势，也增强了开展科研工作的信心。

\noindent\textbf{三、科研训练与课题进展}

\par
\indent
在导师胡晓敏教授的指导下，本人围绕学位论文《基于强化学习惩罚因子的多智能体协同优化方法》持续开展研究工作。课题以网络化分布式优化（Network-based Distributed Optimization，NDO）为研究背景，针对现有多智能体协同优化算法中惩罚参数依赖人工设定、通信协商成本较高以及共享变量处理效率不足等问题，探索将强化学习引入协同优化过程，提高算法在复杂黑箱优化场景下的自适应能力与协同效率。

\par
\indent
围绕上述研究目标，本人完成了强化学习惩罚因子调控模块、共享变量重要性评估机制以及智能通信门控策略等关键模块的设计与实现，并在原有 MACPO 框架基础上构建了 RL-MACPO 算法。通过不断修改程序、分析实验结果和优化算法结构，逐步形成了较完整的算法框架，并完成了核心代码开发、实验平台搭建以及数据统计分析等工作。

\par
\indent
目前，课题已取得较为扎实的阶段性成果。在标准基准测试中，完成了 RL-MACPO 与原始 MACPO 的系统对比实验，并围绕算法关键模块开展了消融实验、强化学习策略对比实验以及通信策略分析。同时，在多个扩展测试函数及典型应用场景中，对算法性能进行了进一步验证，并持续完善相关基线算法复现工作，为后续论文分析提供更加充分的实验依据。

\par
\indent
经过这一阶段的科研训练，本人逐步完成了从文献调研、问题分析、算法设计、程序实现到实验验证和论文撰写的完整科研过程。通过不断发现问题、分析问题和解决问题，不仅提高了科研实践能力，也增强了独立开展科研工作的能力。目前，课题已由算法设计阶段逐步进入实验完善、理论分析和论文撰写阶段，为后续完成学位论文奠定了较好的基础。

\end{titlepage}

% ===== 第 2 页（续；边框与第 1 页相同，顶距对齐首页正文）=====
\begin{titlepage}
    \makeatletter
    \AddToShipoutPictureBG*{
        \begin{tikzpicture}[overlay,remember picture]
            \draw [line width=1pt]
            ($ (current page.north west) + (1.95cm,-2.1cm) $)
            rectangle
            ($ (current page.south east) + (-1.95cm,1.75cm) $);
        \end{tikzpicture}
    }
    \makeatother
    \vspace{2.15cm}

\par
\textbf{(续)}
\par

\textbf{四、主要收获与体会}

\par
\indent
经过硕士阶段的学习和科研训练，本人最大的收获不仅是专业知识的积累，更重要的是科研思维和解决问题能力的提升。刚接触课题时，对网络化分布式优化、多智能体协同进化以及强化学习等内容的理解还停留在理论层面。随着文献阅读的不断深入、实验的反复验证以及导师的耐心指导，我逐渐能够围绕实际问题分析算法设计思路，并尝试提出自己的改进方案。

\par
\indent
科研过程中，我深刻体会到，一个可靠的实验结果往往需要经过大量重复验证才能真正说明问题。从程序调试、参数选择，到实验复现、统计分析，每一个细节都会影响最终结论。正是在一次次修改程序、排查问题和不断优化实验的过程中，我逐渐养成了严谨求实的科研态度，也更加理解了科研工作需要耐心、坚持和不断积累。

\par
\indent
作为非全日制研究生，在兼顾工作与学习的过程中，我更加珍惜每一次科研训练和学习机会，也不断提高了时间管理和自我规划能力。虽然科研过程中遇到过程序调试困难、实验周期较长以及论文写作压力等问题，但每解决一个问题，都让我对科研工作有了更加深入的理解，也更加坚定了继续深入研究的信心。
\noindent\textbf{五、存在不足与改进措施}

\par
\indent
对照培养目标和学位论文要求，本人认识到目前仍存在一些不足。首先，部分扩展实验规模较大、运行周期较长，实验数据仍需进一步补充和完善，后续需要进一步提高实验组织和资源调度效率。其次，论文理论分析部分仍有进一步深化空间，对强化学习调控机制、通信策略以及算法性能提升原因的分析还需更加系统和深入。此外，在论文整体结构和学术表达方面仍需进一步打磨，使论文逻辑更加严谨、内容更加完整。

\par
\indent
下一阶段，本人将在导师指导下继续完善各项实验内容，进一步补充基线算法对比与扩展实验，加强算法机理分析，认真完成学位论文撰写工作，不断提高论文质量。同时，将继续保持严谨务实的科研态度，认真完成硕士阶段各项培养任务，努力以更加扎实的研究成果完成学位论文。

\end{titlepage}
